\section{几段并不平行的年序：南宋人怎样面对国家的变化}

\subsection{建炎：驻跸不是迁都的同义词}

建炎元年以后，朝廷所称的驻跸并非一项礼仪上的临时安排。赵构和随行官员停在扬州、镇江、杭州等处时，必须知道哪里还有粮仓，哪里的渡口可用，哪些军队愿意勤王，哪些州县还能够收取和转运财物。选择一个停驻点，会让那里突然面对官员、卫队、文书、马匹和难民的需求；没有接上这些需求，皇帝所在的地方也难以成为可以发号施令的行在。

苗傅、刘正彦的兵变把这个事实暴露得很尖锐。宫廷不是与城市军队隔绝的中心，而是依靠禁军、地方将领、道路和政治消息才得以存在的脆弱节点。后来关于兵变的叙述往往强调人物忠逆；这些判断不能替代更具体的问题：谁控制宫门和舟船，谁能传递诏令，救援者怎样穿过水网，地方如何看待不断变化的命令。材料允许我们确认政局的急转，不能复原每个参与者的内心和每条街巷的情形。

金军南下与北撤使长江的双重性质更明显。它既是阻碍进攻的水面，也是运送敌我人马、粮食和消息的道路。渡江失败或北撤从来不等于江南天然安全；它可能由补给、季节、抵抗、情报或北方战略共同造成。对南宋而言，保存江南核心地区并没有结束危机，而是把重建的重心转向如何使这些水道持续服务军政。

\subsection{绍兴：议和文本与边城生活}

绍兴和议建立后，使者使用的礼仪称谓、岁币安排和边界用语很快进入史书。它们很重要，因为政治关系必须借具体文本和往返程序才可被承认；但它们没有把边城变成平静的背景。守军仍需粮饷，渡口仍有盘查，沿淮居民仍要面对耕作能否持续、市场是否开放、战事会否再起。所谓“和”，在这里是由哨所、仓场、船队与地方忍耐不断维持的状态。

后世争论和议常将视线固定在朝廷大臣的主战或主和。对边区而言，两种政策都意味着不同的资源配置。远征要求把更多人、粮和船送到前方；守势要求长期维持城寨、巡逻和转运。一个村落或市镇没有义务按照临安的政治语言来感受这一选择。它可能希望道路开放，也可能怕商路引来敌军；它可能靠驻军交易维生，也可能被军需挤压。材料少有机会完整记录这些矛盾，叙事就不应替地方给出整齐的态度。

\subsection{孝宗：恢复的欲望与符离的限制}

高宗禅位后，孝宗即位为恢复议题带来新的政治动力。军事行动的准备却不能只靠君臣意愿。边防部队的驻地、将领之间的协同、军粮能否到达、对金情报是否可靠，以及江南财政能否承担更高强度的支付，都会限制“北伐”究竟能走多远。符离受挫后，历史很容易寻找一个应负责的人或一个应避免的决定；更有解释力的是看清一次远征如何暴露既有补给和组织的薄弱处。

隆兴和议的达成并不是恢复愿望的消散，而是承认某些条件暂时无法由军事解决。它让朝廷重新安排军政和财政，也使边区继续生活在可变的边界中。制度的收益在于减少一次全面碰撞的风险，代价在于把大量军费和外交成本延长到日常。这样的成本往往不以戏剧性的方式出现：它体现在一批粮的转运、一座桥的维修、一条船的征用和一次差役的分配里。

\subsection{宁宗：党禁与前线并非两条历史}

庆元党禁常被放进思想史，开禧北伐则常被放进军事史。两者在当时的国家中并不完全分离。谁能进入朝廷、哪些师友关系被视为可疑、地方书院怎样维持声望，都会影响中枢可用的人和政治信任；而战事准备需要的财政、任命和地方合作又使这些信任成为实际资源。不能因此把党禁解释成北伐的单一原因，或把北伐解释成学术争论的结果；它们只是共同显示政治秩序怎样穿过官场、书院和地方网络。

开禧北伐的挫败尤其警告读者不要把道德判断误作因果解释。韩侂胄的权力、前线指挥、金方反应、兵粮调度与朝中不同意见彼此牵连。战后杀一人、换一约，能解决部分政治责任的归属，却不能使已经付出的运输、征发和市场代价消失。嘉定和议后的边境仍然需要守备，地方仍然要为下一轮不确定性准备粮食、船只与人手。

\subsection{端平：联盟之后的城、路和粮区}

宋蒙共同攻金时，双方可以在军事上暂时指向同一敌人，却不必拥有同一幅战后地图。蔡州陷落后，南宋试图进入河南故地，面对的并非一块等待接收的空地，而是被战争、迁徙和不同军事控制改变过的城池、道路和粮区。旧地的象征意义强烈，实际治理却要求持续供应、安置居民并与新的邻国划定可接受的边界。

宋蒙关系迅速转为冲突，说明联盟的结束不只是外交文书的改变。原先用于进攻的道路、渡口和仓储变成需要防守的对象；地方人要判断接近哪一方能保全家庭和市场；中枢则发现恢复旧疆并不能自动恢复旧有资源。把这一段简写成策略错误，会忽略每一方在信息有限、控制不稳和补给紧张中作出的选择。

\subsection{四川：山城并非孤立的堡垒}

蒙古攻宋后，四川的山城体系为南宋提供了难以迅速跨越的纵深。山城的价值来自地势，却不能只归功于地势：城内要有水、粮和人，城外要有可联络的区域，山道和江流要允许消息与补给在危险中移动。钓鱼城的战事之所以在各种叙述里占据位置，正因为它使这些条件变得格外集中。

蒙哥的死亡有不同传说，可靠材料不能将每一种后出细节都写成事实。无论具体死因如何，这一时段的战事确实改变了蒙古战争的节奏。历史解释应停在材料能承受之处：山城与峡江使征服成本升高，宋方的防御又依赖有限而脆弱的地区供给。把困难归给地理宿命，或把坚持归给单一英雄，都会看不见这些相互条件。

\subsection{襄樊：封锁如何进入江南}

襄阳、樊城的长期围困让战争的时间变得可见。封锁者要保持营地、船队、攻城器械和后方补给；守城者要使汉水、江陵和长江方向的通道保持片刻可用。每一个季节都会影响粮食、疾病、河流与船只。围城持续多年，意味着它不只是一场军人之间的较量，也不断改变周边农田、渡口、商路和征发的节奏。

襄樊失守后，元军获得了通向长江中游的更有利条件，南宋失去了吸收北方冲击的重要关节。这个变化不能被简化成“江南富而不能战”。江南资源确实仍可被征收和调度，但资源必须沿可以通行的道路、经过仍能运作的仓场和港口，才能变成前线所需的粮与器械。空间秩序一旦破裂，原本有效的财政和运输组织便不能按旧方式工作。

\subsection{德祐与祥兴：政权结束后的海路}

临安降元后，朝廷的文书、君主和部分官员进入新安排；南方的行动并未由此停止。宗室、将领、船户和沿海社区沿着福建、广东和海上的路线继续移动，面对的是淡水、粮食、风季、船只与新旧权力的接触。有人选择降附，有人继续抵抗，有人只想保全家人或生意。宋方忠义叙事和元方统一叙事各自把这些人分类，均不能代替复杂而地方化的去留。

崖山之后，南宋政权的年序终止。海域与沿岸社会却继续处理此前留下的问题：船只和港口归谁管理，旧有财产如何被重新登记，迁徙者能否找到落脚点，寺院和家族如何安置死者与记忆。把这一时刻放在较长的地方时间中，并不是削弱政权覆亡的重量，而是承认国家的终结会以不同速度进入家庭、市场和宗教生活。

\subsection{把年序接回地方}

这些片段从建炎到祥兴，显示南宋的历史不能只由一条政治线索承担。高宗的行在依赖渡口与仓储；和议依赖边城、使节和军费；会子依赖市场信任；北伐依赖道路与地方动员；襄樊的围城依赖汉水和江南供给；崖山依赖海上补给与沿岸网络。每一处节点都把国家选择转化为具体的人、物和空间关系。

史料的差异也正是在这些节点上最有意义。《宋史》提供王朝叙述，《建炎以来系年要录》提供高宗朝密集的中枢年序，《宋会要辑稿》保存制度语言；契约、碑刻、寺院材料、港口遗址和城址则让若干地方条件浮现。没有任何一种材料能独自讲完南宋。让它们彼此限制，既能避免把政权命令当作实际生活，也能避免用一处地方的偶然保存重写整个时代。
